% =====================================================================
%  DRAFT USULAN PENELITIAN TUGAS AKHIR
%  Hilma Amalia Putri — Informatika, Fakultas Teknik, Universitas Siliwangi
%
%  Format mengikuti Pedoman Pelaksanaan TA (FT.IF-AKD-F.A.05):
%    - Times New Roman 12 pt, spasi 2
%    - margin kiri 4 cm, atas 4 cm, kanan 3 cm, bawah 3 cm
%    - rata kiri-kanan, penomoran bab angka Romawi
%    - nomor halaman isi berpola I-1, II-1, ...
%
%  Kompilasi: bash build.sh   (lihat README.md)
% =====================================================================

\documentclass[12pt,a4paper,oneside]{report}

% ---------- Mesin & font ---------------------------------------------
% Prioritas: Times New Roman asli bila fonts/times.ttf tersedia.
% Bila tidak, dipakai mathptmx (klon Times) supaya tetap bisa dikompilasi.
\ifdefined\XeTeXversion
  \usepackage{fontspec}
  \IfFileExists{fonts/times.ttf}{
    \setmainfont{times.ttf}[
      Path           = fonts/,
      BoldFont       = timesbd.ttf,
      ItalicFont     = timesi.ttf,
      BoldItalicFont = timesbi.ttf
    ]
  }{
    \usepackage{mathptmx}
  }
\else
  \usepackage[T1]{fontenc}
  \usepackage[utf8]{inputenc}
  \usepackage{mathptmx}
\fi

% ---------- Bahasa ----------------------------------------------------
\usepackage[english,indonesian]{babel}

% ---------- Tata letak halaman (pedoman 2.2.1) ------------------------
\usepackage[a4paper,top=4cm,left=4cm,bottom=3cm,right=3cm]{geometry}
\usepackage{setspace}
\doublespacing                 % spasi 2 sesuai pedoman
\usepackage{indentfirst}       % paragraf pertama ikut menjorok
\setlength{\parindent}{1.27cm} % 7 ketukan / 1 Tab

% ---------- Nomor halaman: I-1, II-1, ... -----------------------------
% Gaya halaman buatan sendiri agar nomor halaman hanya muncul di kanan atas
% pada halaman isi. Halaman pertama setiap bab memakai gaya plain bawaan
% (nomor di bawah tengah) sesuai pedoman 2.2.2d.
\makeatletter
\def\ps@isita{%
  \let\@mkboth\@gobbletwo
  \def\@oddhead{\hfil\normalfont\normalsize\thepage}%
  \def\@evenhead{\normalfont\normalsize\thepage\hfil}%
  \def\@oddfoot{}\def\@evenfoot{}}
\makeatother
\pagestyle{isita}
\renewcommand{\thepage}{\Roman{chapter}-\arabic{page}}

% Setiap bab dimulai lagi dari halaman 1. Reset dilakukan dengan etoolbox
% agar tidak merusak \chapter* (misalnya ABSTRAK dan KATA PENGANTAR).
\usepackage{etoolbox}
\pretocmd{\chapter}{\setcounter{page}{1}}{}{}

% Lebar kolom nomor halaman pada daftar isi (agar "III-4" tidak melewati batas)
\makeatletter
\renewcommand{\@pnumwidth}{2.4em}
\renewcommand{\@tocrmarg}{3.2em}
\makeatother

% ---------- Judul bab, subbab, subsubbab (pedoman 2.2.3) --------------
\usepackage{titlesec}
\titleformat{\chapter}[display]
  {\centering\bfseries\normalsize}
  {BAB \Roman{chapter}}{0pt}{\MakeUppercase}
\titlespacing*{\chapter}{0pt}{0pt}{1.5\baselineskip}
\titleformat{\section}{\bfseries\normalsize}{\thesection}{0.6em}{}
\titleformat{\subsection}{\bfseries\normalsize}{\thesubsection}{0.6em}{}
\titleformat{\subsubsection}{\bfseries\normalsize}{\thesubsubsection}{0.6em}{}

% ---------- Tabel, gambar, kode ---------------------------------------
\usepackage{graphicx}
\usepackage{float}
\usepackage{booktabs}
\usepackage{tabularx}
\usepackage{array}
\usepackage{longtable}
\usepackage{caption}
\captionsetup{
  font            = small,
  labelfont       = bf,
  labelsep        = period,
  justification   = centering,
  singlelinecheck = false
}
\captionsetup[table]{position=top, skip=0.5\baselineskip}
\captionsetup[figure]{position=bottom, skip=0.5\baselineskip}
\usepackage{etoolbox}
\AtBeginEnvironment{longtable}{\singlespacing\small}
\AtBeginEnvironment{tabularx}{\singlespacing\small}

% Penomoran tabel dan gambar mengikuti bab
\counterwithin{table}{chapter}
\counterwithin{figure}{chapter}
\renewcommand{\thetable}{\Roman{chapter}.\arabic{table}}
\renewcommand{\thefigure}{\Roman{chapter}.\arabic{figure}}

% ---------- Daftar poin: huruf/angka, bukan bullets (pedoman 2.2.1) ----
\usepackage{enumitem}
\newlist{poin}{enumerate}{2}
\setlist[poin]{label=\alph*., ref=\alph*., leftmargin=*, align=left,
  itemsep=0pt, topsep=0.25\baselineskip}
\newlist{urut}{enumerate}{2}
\setlist[urut]{label=\arabic*., ref=\arabic*., leftmargin=*, align=left,
  itemsep=0pt, topsep=0.25\baselineskip}

% ---------- Matematika & diagram --------------------------------------
\usepackage{amsmath}
\usepackage{tikz}
\usetikzlibrary{shapes.geometric,arrows.meta,positioning,calc}

% ---------- Sitasi: natbib + BibTeX, gaya (Nama, Tahun) ---------------
\usepackage[round,authoryear]{natbib}
\bibliographystyle{plainnat}
\newcommand{\parencite}[1]{\citep{#1}}
\newcommand{\textcite}[1]{\citet{#1}}

% ---------- Tautan (dimatikan agar tidak berwarna saat dicetak) -------
\usepackage[hidelinks]{hyperref}

% ---------- Data dokumen ----------------------------------------------
\newcommand{\JudulTA}{VERIFIKASI BERLAPIS UNTUK MENEKAN KESALAHAN SENYAP
  PADA TAMPILAN TUJUH SEGMEN TIMBANGAN DIGITAL}
\newcommand{\NamaPenulis}{Hilma Amalia Putri}
\newcommand{\NPMPenulis}{[NPM]}
\newcommand{\NamaPembimbingSatu}{[Nama Pembimbing I]}
\newcommand{\NamaPembimbingDua}{[Nama Pembimbing II]}
\newcommand{\BulanTahun}{[Bulan] 2026}

% Perintah penanda bagian yang belum diisi
\newcommand{\isian}[1]{{\itshape[belum diisi: #1]}}

% =====================================================================
\begin{document}

% ---------------- Halaman sampul (Lampiran 1 pedoman) ----------------
\begin{titlepage}
  \centering
  \vspace*{1cm}
  {\bfseries\normalsize \JudulTA \par}
  \vspace{2cm}
  {\bfseries\normalsize USULAN PENELITIAN\\TUGAS AKHIR\par}
  \vspace{2cm}
  {\normalsize Oleh:\par}
  \vspace{0.5cm}
  {\normalsize Nama: \NamaPenulis\par}
  {\normalsize NPM: \NPMPenulis\par}
  \vspace{1.5cm}
  {\bfseries\normalsize JURUSAN INFORMATIKA\\[0.2cm]
   FAKULTAS TEKNIK UNIVERSITAS SILIWANGI\\[0.2cm]
   TASIKMALAYA\par}
  \vspace{1cm}
  {\normalsize \BulanTahun\par}
\end{titlepage}

% ---------------- Bagian awal: angka Romawi kecil --------------------
\pagenumbering{roman}
\pagestyle{plain}

\chapter*{ABSTRAK}
\addcontentsline{toc}{chapter}{ABSTRAK}
\singlespacing
% =====================================================================
%  ABSTRAK — ditulis dalam satu paragraf, memuat masalah, metode,
%  dan kesimpulan. Angka hasil diisi setelah eksperimen selesai.
% =====================================================================

Pencatatan berat barang di gudang dan pabrik masih dilakukan secara manual,
sehingga kesalahan catat tidak terdeteksi dan tidak tersedia jejak audit.
Pengenalan karakter optis untuk tampilan tujuh segmen telah berkembang, tetapi
seluruh penelitian yang ada berhenti pada pelaporan akurasi dan tidak
melaporkan bacaan salah yang lolos tanpa peringatan. Penelitian ini merancang
sistem pembacaan tampilan tujuh segmen dengan empat lapisan verifikasi
kontekstual, yaitu verifikasi struktural, fisik, temporal, dan statistik, di
atas keluaran PP-OCRv6 yang dijalankan tanpa koneksi jaringan. Data latih
dibangkitkan secara sintetis dengan kondisi terkontrol, sedangkan data uji
berupa citra timbangan nyata yang dikunci sebelum eksperimen. Pengujian dilakukan
berjenjang melalui enam eksperimen untuk memisahkan kontribusi setiap lapisan.
Keberhasilan diukur dengan laju kesalahan senyap, \textit{coverage},
\textit{selective accuracy}, dan luas daerah di bawah kurva
\textit{risk--coverage}. Hasil penelitian menunjukkan bahwa sistem menurunkan
laju kesalahan senyap dari \isian{angka E1} menjadi \isian{angka E5} pada
\textit{coverage} sebesar \isian{angka}, dengan waktu pemrosesan
\isian{angka} milidetik per transaksi. Dengan demikian, penelitian ini
memberikan kerangka evaluasi keandalan yang belum tersedia pada literatur
pembacaan tampilan tujuh segmen.

\vspace{1\baselineskip}
\noindent\textbf{Kata kunci:} tampilan tujuh segmen, kesalahan senyap, verifikasi
kontekstual, PP-OCRv6, \textit{selective prediction}

\doublespacing

\chapter*{KATA PENGANTAR}
\addcontentsline{toc}{chapter}{KATA PENGANTAR}
% =====================================================================
%  KATA PENGANTAR
% =====================================================================

Puji syukur penulis panjatkan kepada Allah Swt. atas rahmat dan karunia-Nya,
sehingga usulan penelitian tugas akhir yang berjudul
``\JudulTA'' ini dapat disusun.

Usulan penelitian ini disusun sebagai salah satu syarat dalam menempuh mata
kuliah Tugas Akhir pada Jurusan Informatika, Fakultas Teknik, Universitas
Siliwangi. Penelitian ini membahas perancangan sistem pembacaan tampilan tujuh
segmen pada timbangan digital yang dilengkapi lapisan verifikasi kontekstual,
dengan penekanan pada pengukuran keandalan dan bukan semata-mata akurasi.

Pada kesempatan ini penulis menyampaikan terima kasih kepada:
\begin{poin}
  \item \NamaPembimbingSatu, selaku pembimbing I, atas arahan mengenai
        substansi dan arah penelitian;
  \item \NamaPembimbingDua, selaku pembimbing II, atas arahan mengenai tata
        tulis dan teknis pengerjaan;
  \item pihak Jurusan Informatika yang telah memberikan kemudahan administratif
        selama penyusunan usulan penelitian ini.
\end{poin}

Penulis menyadari bahwa usulan penelitian ini masih memiliki kekurangan. Oleh
karena itu, kritik dan saran yang membangun sangat diharapkan.

\vspace{1\baselineskip}
\begin{flushright}
Tasikmalaya, \BulanTahun\\
Penyusun
\end{flushright}


\clearpage
\tableofcontents
\clearpage
\listoftables
\clearpage
\listoffigures

% ---------------- Bagian isi: I-1, II-1, ... -------------------------
\clearpage
\pagenumbering{arabic}
\renewcommand{\thepage}{\Roman{chapter}-\arabic{page}}
\pagestyle{isita}

% =====================================================================
%  BAB I PENDAHULUAN
% =====================================================================

\chapter{PENDAHULUAN}

\section{Latar Belakang}

Pencatatan berat barang di gudang dan pabrik masih banyak dilakukan secara
manual. Operator membaca angka pada tampilan timbangan digital, kemudian
menuliskannya ke buku catatan atau lembar kerja. Cara tersebut memiliki dua
kelemahan yang melekat. Pertama, tidak tersedia pembanding otomatis, sehingga
kesalahan catat tidak dapat dideteksi pada saat kejadian. Kedua, tidak
tertinggal jejak audit, sehingga sengketa jumlah barang sulit ditelusuri.
Kebutuhan akan perekaman otomatis telah dijawab di literatur, antara lain
melalui sistem pencatat data waktu nyata berbasis \textit{Convolutional
Recurrent Neural Network} dan \textit{Connectionist Temporal Classification}
(CRNN-CTC) untuk tampilan tujuh segmen \parencite{effendi2025unjuk}. Karena
tampilan timbangan sudah menyajikan angka secara digital, pembacaan sebenarnya
dapat diotomasi menggunakan kamera dan pengenalan karakter optis tanpa perlu
mengganti perangkat timbangan yang telah terpasang.

Persoalannya, mesin pengenalan karakter yang dirancang untuk teks dokumen tidak
bekerja dengan baik pada tampilan tujuh segmen. \textcite{low2024detecting}
melaporkan bahwa pengenalan pada kumpulan citra tujuh segmen tetap bermasalah,
terutama pada kondisi pantulan cahaya. Penyebabnya bersifat struktural:
karakter pada tampilan tujuh segmen bukan huruf, melainkan kombinasi segmen yang
menyala. Angka ``1'' hanya berupa dua garis tegak, sedangkan angka ``8''
menyalakan ketujuh segmen. Jarak antardigit sangat rapat, dan titik desimal
berukuran hanya beberapa piksel sehingga sering tidak terbaca sebagai karakter
tersendiri. Persoalan desimal inilah yang menyebabkan
\textcite{gushima2023automatic} memisahkan model digit dari model titik desimal.

Sejumlah pendekatan khusus kemudian dikembangkan. \textcite{laroca2019convolutional}
membangun metode berbasis jaringan konvolusional dan merilis kumpulan data
UFPR-AMR berisi 2.000 citra meter energi beserta 10.000 anotasi tingkat digit.
\textcite{bonilla2023adapting} mengadaptasi YOLO untuk mengenali angka riil
termasuk titik desimal. \textcite{yenneti2023lightweight} mengusulkan kepala
inferensi digit yang mencapai 97,8 persen pada kerangka tunggal dengan kecepatan
30 bingkai per detik, disertai penyaringan antarbingkai.
\textcite{wei2023novel} melengkapinya dengan model sekuens bermemori pendek
untuk panjang bacaan yang bervariasi. Di sisi lain,
\textcite{jiang2026method} justru menunjukkan bahwa pendekatan klasik berbasis
binarisasi adaptif, proyeksi, dan pencocokan Hamming mencapai 99,7 persen dalam
20 milidetik pada timbangan elektronik. Temuan tersebut menandakan bahwa
pembelajaran mendalam tidak serta-merta mengungguli metode klasik pada domain
yang geometrinya sangat teratur.

Pada periode berikutnya perhatian bergeser ke model yang lebih besar dan lebih
umum. \textcite{ega2026leveraging} membandingkan model bahasa-visi dengan lima
mesin pengenalan karakter pada perangkat ukur digital dan melaporkan bahwa
Gemini mencapai \textit{full-string recognition} sebesar 93,85 persen.
Angka tersebut tampak bertentangan dengan temuan \textcite{lin2026measurebench}
yang melaporkan model bahasa-visi terbaik hanya mencapai 30,2 persen pada data
nyata. Keduanya tidak bertentangan: pengujian pertama menilai transkripsi digit
yang telah tertulis, sedangkan pengujian kedua menuntut inferensi skala dan
penalaran spasial. Perbedaan itu menempatkan model bahasa-visi bukan sebagai
pesaing akurasi pada domain transkripsi digit, melainkan sebagai pesaing biaya.

Pada Juni 2026 dirilis PP-OCRv6, keluarga model pengenalan karakter ringan
berukuran 1,5 hingga 34,5 juta parameter \parencite{zhang2026ppocrv6}. Model
tersebut memakai tulang punggung PPLCNetV4, leher pengenalan
\textit{EncoderWithLightSVTR}, serta dekoder gabungan CTC dan NRTR. Pada tolok
ukur internal pengembangnya, tingkat \textit{medium} mencapai 83,2 persen
akurasi pengenalan rata-rata dari 15 kategori dan secara eksplisit menyebut
tampilan digital sebagai salah satu sasaran khusus
\parencite{zhang2026ppocrv6}. Kehadirannya melanjutkan kecenderungan yang telah
terlihat pada PP-OCRv5, yaitu model berukuran 5 juta parameter yang mampu
menyaingi model bahasa-visi berukuran miliaran parameter pada tugas serupa
\parencite{zhang2026ppocrv5}.

Seluruh pekerjaan di atas berhenti pada titik yang sama: melaporkan akurasi pada
kumpulan data yang disusun sendiri. Angka yang beredar tidak dapat dibandingkan
satu sama lain karena dihasilkan dari data, pembagian latih-uji, dan definisi
``benar'' yang berbeda. Persoalan yang lebih mendasar terletak pada apa yang
dilaporkan. Semua sumber tersebut mengukur akurasi, tetapi tidak mengukur
kesalahan yang paling mahal dalam pencatatan, yaitu bacaan yang salah namun
lolos tanpa peringatan. Pada konteks penimbangan, satu bacaan salah lebih mahal
dampaknya daripada sepuluh bacaan yang ditolak, karena bacaan salah masuk ke
catatan dan baru diketahui setelah transaksi selesai. Sistem nyata memang telah
memakai ambang keyakinan, tetapi ambang tersebut dipilih tanpa divalidasi
terhadap data, padahal nilai keyakinan model bukan bukti kebenaran.
Sepanjang penelusuran penulis terhadap 24 sumber terbitan 2019--2026, tidak
ditemukan penelitian yang memakai konteks di luar citra untuk memverifikasi
hasil pembacaan tampilan tujuh segmen.

Berdasarkan uraian tersebut, penelitian ini mengambil tiga hal yang belum
dikerjakan sumber sebelumnya. Pertama, menguji kemampuan PP-OCRv6 dalam membaca
tampilan tujuh segmen tanpa adaptasi dan menentukan adaptasi minimum yang
diperlukan \parencite{zhang2026ppocrv6}. Kedua, merancang lapisan verifikasi
kontekstual berlapis yang memanfaatkan konteks di luar citra untuk memutuskan
apakah suatu bacaan layak dipercaya. Ketiga, mengukur keandalan sistem tidak
hanya melalui akurasi, tetapi juga melalui laju kesalahan senyap,
\textit{coverage}, \textit{selective accuracy}, dan luas daerah di bawah kurva
\textit{risk--coverage}.

\section{Rumusan Masalah}

Berdasarkan latar belakang tersebut, rumusan masalah penelitian ini adalah
sebagai berikut.
\begin{poin}
  \item Sejauh mana PP-OCRv6 mampu membaca tampilan tujuh segmen tanpa
        adaptasi, dan adaptasi minimum apa yang diperlukan agar kemampuannya
        memadai?
  \item Bagaimana merancang lapisan verifikasi kontekstual yang dapat menekan
        laju kesalahan senyap pada pembacaan tampilan tujuh segmen?
  \item Berapa \textit{trade-off} antara \textit{coverage} dan akurasi pada
        bacaan yang diterima, serta bagaimana dampaknya terhadap beban koreksi
        operator?
\end{poin}

\section{Tujuan Penelitian}

Tujuan yang ingin dicapai penelitian ini adalah sebagai berikut.
\begin{poin}
  \item Mengukur kemampuan PP-OCRv6 pada tampilan tujuh segmen dalam kondisi
        tanpa adaptasi dan menentukan tingkat adaptasi minimum yang diperlukan.
  \item Merancang dan mengimplementasikan empat lapis verifikasi kontekstual,
        yaitu verifikasi struktural, fisik, temporal, dan statistik.
  \item Menghasilkan bukti kuantitatif keandalan sistem berupa laju kesalahan
        senyap, \textit{coverage}, \textit{selective accuracy}, dan luas daerah
        di bawah kurva \textit{risk--coverage}, beserta dampaknya terhadap
        tingkat koreksi operator.
\end{poin}

\section{Manfaat Penelitian}

Penelitian ini diharapkan memberikan manfaat sebagai berikut.
\begin{poin}
  \item Bagi pemilik usaha dan operator, tersedia catatan berat yang dapat
        diaudit, dan bacaan yang meragukan tidak lagi masuk ke catatan secara
        diam-diam.
  \item Bagi pengembangan ilmu, tersedia kerangka evaluasi keandalan untuk
        domain tampilan tujuh segmen, yang selama ini hanya diukur dengan
        akurasi \parencite{effendi2025unjuk,gushima2023automatic,yenneti2023lightweight}.
  \item Bagi praktik rekayasa, tersedia bukti kuantitatif mengenai biaya
        penolakan, yaitu berapa banyak bacaan yang harus dikorbankan untuk
        mencapai tingkat kepercayaan tertentu, sehingga pemilihan ambang dapat
        dilakukan berdasarkan data dan bukan perkiraan.
\end{poin}

\section{Batasan Masalah}

Agar pembahasan terarah, penelitian ini dibatasi pada hal-hal berikut.
\begin{poin}
  \item Karakter yang dikenali terbatas pada angka 0--9 dan titik desimal,
        dengan panjang 1--8 karakter.
  \item Objek penelitian adalah tampilan tujuh segmen jenis \textit{light
        emitting diode} dan \textit{liquid crystal display} pasif pada timbangan
        digital.
  \item Data uji berupa citra nyata yang tidak dipakai untuk pelatihan maupun
        penyetelan, dan dikunci sebelum eksperimen dijalankan.
  \item Data latih dihasilkan secara sintetis dengan kondisi terkontrol.
  \item Sistem dirancang berjalan mandiri di perangkat tanpa koneksi jaringan,
        dengan sinkronisasi data dilakukan saat jaringan tersedia.
  \item Lingkup sistem terbatas pada satu lini produk tanpa integrasi dengan
        sistem perencanaan sumber daya perusahaan.
\end{poin}

% =====================================================================
%  BAB II LANDASAN TEORI
% =====================================================================

\chapter{LANDASAN TEORI}

\section{Tampilan Tujuh Segmen}

Tampilan tujuh segmen menyajikan angka menggunakan tujuh batang segmen yang
disusun seperti angka delapan, ditambah satu titik desimal. Setiap digit
dibentuk oleh kombinasi segmen yang menyala, bukan oleh bentuk huruf. Karena
itu pendekatan pengenalan berbasis huruf tidak dapat langsung diterapkan.
Kekhasan tersebut terlihat pada angka ``1'' yang hanya menyalakan dua segmen
tegak dan angka ``8'' yang menyalakan seluruh segmen, sehingga jarak kesamaan
antarkelas tidak seragam. Titik desimal berukuran sangat kecil, umumnya tiga
hingga lima piksel pada citra beresolusi sedang, sehingga rawan hilang pada
proses pengenalan \parencite{gushima2023automatic}.

Tampilan dibedakan menjadi dua jenis berdasarkan sumber cahayanya. Tampilan
\textit{light emitting diode} memancarkan cahaya sendiri, umumnya berwarna
merah, hijau, atau biru, dengan latar gelap. Tampilan \textit{liquid crystal
display} pasif tidak memancarkan cahaya dan bergantung pada cahaya sekitar,
sehingga kontrasnya rendah dan sudut pandangnya terbatas. Kedua jenis tersebut
memerlukan penanganan praproses yang berbeda. Koreksi pencahayaan pada kondisi
minim cahaya yang tidak seragam telah diteliti oleh \textcite{fu2026twolevel}.

\section{Pengenalan Karakter Optis untuk Deret Angka}

Pengenalan deret angka tanpa segmentasi karakter umumnya memakai arsitektur
CRNN yang diperkenalkan \textcite{shi2017crnn}. Arsitektur tersebut menggabungkan
lapisan konvolusional untuk mengekstraksi ciri, lapisan \textit{bidirectional
long short-term memory} untuk memodelkan urutan, dan fungsi kerugian CTC untuk
menyelaraskan urutan keluaran dengan label tanpa segmentasi eksplisit
\parencite{graves2006ctc}.

CTC memiliki keterbatasan struktural yang relevan bagi penelitian ini. Karena
segmentasi karakter dianggap tidak diketahui, dua karakter identik yang
berurutan hanya dapat dihasilkan apabila terdapat simbol kosong di antaranya.
Akibatnya deret seperti ``1006.5'' berpotensi terbaca ``106.5''. Pola kegagalan
tersebut terdokumentasi pada implementasi CRNN untuk tampilan tujuh segmen
\parencite{oicws2026lcd}. Alternatif dekoder tanpa rekurensi dikembangkan
melalui NRTR \parencite{sheng2019nrtr} dan PARSeq
\parencite{bautista2022parseq}.

\section{PP-OCRv6}

PP-OCRv6 adalah keluarga model pengenalan karakter yang dirilis pada Juni 2026
\parencite{zhang2026ppocrv6}. Model ini tersedia dalam tiga tingkat, yaitu
\textit{tiny} berukuran 1,5 juta parameter, \textit{small} berukuran 7,7 juta
parameter, dan \textit{medium} berukuran 34,5 juta parameter. Pada tingkat
\textit{small}, kepala pengenalannya berukuran sekitar 5,2 juta parameter,
sehingga angka tersebut tidak sama dengan ukuran keseluruhan model.

Secara arsitektural, PP-OCRv6 memakai tulang punggung PPLCNetV4 yang bergaya
\textit{MetaFormer} dengan reparameterisasi struktural, leher deteksi RepLKFPN,
serta leher pengenalan \textit{EncoderWithLightSVTR} yang memakai atensi
lokal-global. Pengenalan dilakukan oleh dekoder gabungan CTC dan NRTR, dengan
hanya cabang CTC yang dievaluasi saat inferensi; cabang NRTR berperan pada saat
pelatihan sebagai pengatur tambahan bagi representasi enkoder
\parencite{zhang2026ppocrv6}. Model ini tersedia melalui Paddle, Transformers,
dan ONNX Runtime, sehingga dapat dijalankan tanpa koneksi jaringan apabila
bobotnya dibundel bersama aplikasi. Model generasi sebelumnya, PP-OCRv5,
menunjukkan bahwa model berukuran 5 juta parameter dapat menyaingi model
bahasa-visi berukuran miliaran parameter pada tugas pengenalan karakter
\parencite{zhang2026ppocrv5}.

\section{Model Bahasa-Visi}

Model bahasa-visi menerima citra dan pertanyaan dalam bentuk teks, lalu
menghasilkan jawaban. Pada domain alat ukur, kemampuan model ini terbelah.
\textcite{ega2026leveraging} melaporkan \textit{full-string recognition} sebesar
93,85 persen dan akurasi per digit 98,10 persen untuk Gemini pada perangkat ukur
digital. Sebaliknya, \textcite{lin2026measurebench} melaporkan bahwa model
terbaik hanya mencapai 30,2 persen pada data nyata dari 26 jenis alat ukur.
Perbedaan tersebut dijelaskan oleh jenis tugasnya: transkripsi digit yang telah
tertulis versus inferensi nilai dari posisi jarum dan skala. Pada domain
penelitian ini, yaitu transkripsi digit, model bahasa-visi bukan pesaing
akurasi, melainkan pesaing dari sisi biaya komputasi dan kebutuhan jaringan.

\section{Keandalan dan \textit{Selective Prediction}}

Akurasi dan keandalan merupakan dua besaran yang berbeda. Akurasi mengukur
seberapa sering model benar, sedangkan keandalan menyangkut kemampuan sistem
menyatakan bahwa dirinya tidak layak dipercaya pada suatu kasus.
\textit{Selective prediction} memberi sistem kewenangan menolak memprediksi
apabila keyakinannya rendah \parencite{lin2026measurebench}. Pada kerangka ini,
kinerja tidak dinyatakan sebagai satu angka, melainkan sebagai hubungan antara
\textit{coverage}, yaitu proporsi kasus yang diterima otomatis, dan akurasi pada
kasus yang diterima. Hubungan tersebut digambarkan sebagai kurva
\textit{risk--coverage} dan diringkas sebagai luas daerah di bawah kurva atau
\textit{area under the risk--coverage curve} (AURC).

Nilai keyakinan mentah dari model tidak dapat langsung dipakai sebagai peluang
kebenaran. Diperlukan kalibrasi pada data yang terpisah dari data penilaian,
agar ambang keputusan tidak ditentukan dari data yang sama dengan data uji.

\section{Penelitian Terdahulu}

Ringkasan penelitian terdahulu yang menjadi acuan disajikan pada
Tabel~\ref{tab:terdahulu}.

\begin{table}[H]
\caption{Ringkasan penelitian terdahulu}
\label{tab:terdahulu}
\begin{tabularx}{\textwidth}{@{}p{0.6cm} p{3.3cm} p{4.6cm} X@{}}
\toprule
\textbf{No} & \textbf{Sumber} & \textbf{Hasil yang dilaporkan} & \textbf{Keterbatasan untuk penelitian ini} \\
\midrule
1 & \textcite{jiang2026method} & 99,7 persen pada 1.146 karakter dalam 20 milidetik, diuji pada timbangan elektronik & Bergantung pada kondisi citra; perilaku saat gagal tidak dilaporkan \\
2 & \textcite{effendi2025unjuk} & \textit{Exact match} 90,37 persen dengan CER 3,6 persen pada data logger waktu nyata & Berhenti pada akurasi; perilaku saat bacaan salah tidak dilaporkan \\
3 & \textcite{gushima2023automatic} & Generator citra sintetis dan model desimal terpisah, 96,8 persen & Terbatas pada perangkat berspesifikasi rendah; keandalan tidak diukur \\
4 & \textcite{yenneti2023lightweight} & 97,8 persen per bingkai pada 30 bingkai per detik, dengan penyaringan antarbingkai & Penyaringan dipakai untuk stabilisasi keluaran, bukan keputusan terima atau tolak \\
5 & \textcite{oicws2026lcd} & 98 persen pada data nyata, bobot terbuka & Angka terikat pada kumpulan data sendiri \\
6 & \textcite{low2024detecting} & DBNet terbaik untuk deteksi, PARSeq terbaik untuk pengenalan & Gagal pada kondisi pantulan cahaya \\
7 & \textcite{ega2026leveraging} & Gemini dengan \textit{full-string recognition} 93,85 persen & Membutuhkan jaringan dan biaya; keandalan tidak diukur \\
8 & \textcite{lin2026measurebench} & Model bahasa-visi terbaik hanya 30,2 persen pada data nyata & Menguji inferensi skala, bukan transkripsi digit \\
9 & \textcite{zhang2026ppocrv6} & PP-OCRv6, 1,5--34,5 juta parameter, 83,2 persen rata-rata 15 kategori & Belum diuji pada tampilan tujuh segmen \\
\bottomrule
\end{tabularx}
\end{table}

\section{Posisi Penelitian}

Berdasarkan kajian pada bagian sebelumnya, terdapat tiga celah yang belum
terjawab. Pertama, keandalan pembacaan tampilan tujuh segmen belum pernah
dilaporkan dalam bentuk laju kesalahan senyap, kurva \textit{risk--coverage},
maupun \textit{selective accuracy}. Kedua, belum ada penelitian yang memakai
konteks di luar citra, seperti jumlah slot digit, rentang berat yang sah untuk
suatu produk, toleransi penimbangan, dan konsistensi antarbingkai, sebagai dasar
keputusan kelayakan bacaan. Ketiga, PP-OCRv6 yang secara eksplisit menyebut
tampilan digital sebagai sasaran khusus belum diuji pada tampilan tujuh segmen
\parencite{zhang2026ppocrv6}.

Penelitian ini menempati celah tersebut dengan menggeser pertanyaan dari sekadar
``apakah bacaan dapat diperoleh'' menjadi ``kapan bacaan tidak layak dipercaya
dan bagaimana sistem mengetahuinya''. Pergeseran tersebut tidak bergantung pada
penggantian model, sehingga kontribusinya tidak dapat direduksi menjadi
pemakaian model yang lebih baru.

% =====================================================================
%  BAB III METODOLOGI
% =====================================================================

\chapter{METODOLOGI}

\section{Rancangan Penelitian}

Penelitian ini merupakan penelitian rekayasa yang membangun sistem pembacaan
tampilan tujuh segmen dengan lapisan verifikasi kontekstual. Alur kerja dibagi
menjadi lima tahap, yaitu:
\begin{urut}
  \item penetapan spesifikasi data dan metrik keberhasilan;
  \item pembangkitan data latih sintetis dan pengumpulan data uji nyata;
  \item pengukuran kemampuan awal PP-OCRv6 tanpa adaptasi;
  \item pembangunan lapisan verifikasi secara bertahap dan terukur;
  \item evaluasi sistem penuh serta analisis dampaknya terhadap operator.
\end{urut}

Setiap lapisan verifikasi dibangun dan diukur secara terpisah. Pemisahan ini
disengaja agar penyebab perubahan laju kesalahan senyap dapat diidentifikasi,
sesuatu yang tidak mungkin dilakukan bila beberapa lapisan diuji sekaligus.

\section{Data dan Instrumen}

\subsection{Data Latih Sintetis}

Data latih dibangkitkan secara terprogram dengan variasi terkontrol, meliputi
jenis font segmen, warna dan intensitas pancaran, ketebalan segmen, kontras
tampilan \textit{liquid crystal display} pasif, sudut kamera, derajat
kemiringan, derau sensor, serta panjang bacaan 1 sampai 8 karakter. Prinsip yang
dipedomani adalah variasi lebih penting daripada jumlah; citra yang banyak namun
seragam tidak menambah kemampuan generalisasi. Pemanfaatan pembangkit sintetis
untuk keperluan serupa telah dilakukan \textcite{gushima2023automatic} dan
\textcite{yenneti2023lightweight}.

Rincian jumlah citra, sebaran kondisi, dan format label ditetapkan pada dokumen
spesifikasi data tersendiri. Data latih tidak pernah dipakai sebagai data
penilaian.

\subsection{Data Uji Nyata}

Data uji berupa citra timbangan digital nyata yang diambil dari sejumlah merek
berbeda pada beberapa kondisi pencahayaan dan sudut pengambilan. Setiap
pembacaan disertai rekaman video pendek agar lapisan verifikasi temporal dapat
diuji. Kumpulan data uji dikunci sebelum eksperimen dijalankan, dan setelah
dikunci tidak boleh dipakai untuk memilih ambang, menyetel praproses, maupun
menentukan hyperparameter.

Ukuran kumpulan data uji dan konsekuensinya terhadap lebar interval kepercayaan
dilaporkan secara eksplisit. Pengumpulan data nyata tetap diperlukan karena
tidak tersedia kumpulan data publik berupa timbangan dengan label tingkat string
beserta titik desimal dan metadata kondisi.

\subsection{Instrumen}

Perangkat lunak yang dipakai adalah PP-OCRv6 \parencite{zhang2026ppocrv6} yang
dijalankan melalui ONNX Runtime agar dapat beroperasi tanpa koneksi jaringan,
serta pustaka pengolahan citra untuk praproses. Sebagai pembanding dipakai
CRNN-CTC \parencite{shi2017crnn} dan perangkat klasik \textit{seven segment
optical character recognition} \parencite{auerswal_ssocr}.

\section{Rancangan Lapisan Verifikasi}

Sistem menerima citra, memotong area tampilan, melakukan praproses, lalu
menjalankan pengenalan. Keluaran pengenalan tidak langsung diterima, melainkan
melewati empat lapisan verifikasi seperti tampak pada Gambar~\ref{fig:alur}.

\begin{figure}[H]
\centering
\begin{tikzpicture}[
  node distance=0.5cm,
  kotak/.style={rectangle, draw, rounded corners=2pt, align=center,
                text width=9.6cm, inner sep=4pt, font=\small},
  lapis/.style={rectangle, draw, thick, rounded corners=2pt, align=center,
                text width=9.6cm, inner sep=4pt, font=\small},
  akhir/.style={rectangle, draw, very thick, rounded corners=2pt,
                align=center, text width=3.6cm, inner sep=5pt, font=\small},
  panah/.style={-{Latex[length=2mm]}, thick}
]
\node[kotak] (aku) {Akuisisi citra atau 3--5 bingkai video};
\node[kotak, below=of aku] (pra) {Pemotongan area tampilan dan praproses};
\node[kotak, below=of pra] (ocr) {PP-OCRv6: deteksi dan pengenalan};
\node[kotak, below=of ocr] (kan) {Kandidat string dan nilai keyakinan};

\node[lapis, below=0.7cm of kan] (l1)
  {\textbf{Lapis 1 --- Verifikasi Struktural}\\[1pt]
   {\scriptsize gagal $\rightarrow$ tolak}};
\node[lapis, below=of l1] (l2)
  {\textbf{Lapis 2 --- Verifikasi Fisik}\\[1pt]
   {\scriptsize gagal $\rightarrow$ tolak}};
\node[lapis, below=of l2] (l3)
  {\textbf{Lapis 3 --- Verifikasi Temporal}\\[1pt]
   {\scriptsize bingkai tidak sepakat $\rightarrow$ minta konfirmasi}};
\node[lapis, below=of l3] (l4)
  {\textbf{Lapis 4 --- Verifikasi Statistik}\\[1pt]
   {\scriptsize keyakinan rendah $\rightarrow$ minta konfirmasi}};

\node[akhir, below=0.7cm of l4] (terima) {\textbf{TERIMA}};

\draw[panah] (aku) -- (pra);
\draw[panah] (pra) -- (ocr);
\draw[panah] (ocr) -- (kan);
\draw[panah] (kan) -- (l1);
\draw[panah] (l1) -- node[right, font=\scriptsize, pos=0.45] {lolos} (l2);
\draw[panah] (l2) -- node[right, font=\scriptsize, pos=0.45] {lolos} (l3);
\draw[panah] (l3) -- node[right, font=\scriptsize, pos=0.45] {lolos} (l4);
\draw[panah] (l4) -- node[right, font=\scriptsize, pos=0.45] {lolos} (terima);
\end{tikzpicture}
\caption{Alur pembacaan dengan empat lapisan verifikasi kontekstual}
\label{fig:alur}
\end{figure}

\subsection{Lapis 1: Verifikasi Struktural}

Lapisan ini memeriksa kesesuaian jumlah digit yang terbaca dengan jumlah slot
yang tersedia pada tampilan, serta kesahihan posisi titik desimal. Pemeriksaan
ini menyasar pola kegagalan yang bersifat struktural, yaitu hilangnya karakter
berulang pada dekoder CTC \parencite{graves2006ctc,oicws2026lcd}.

\subsection{Lapis 2: Verifikasi Fisik}

Lapisan ini memeriksa kesesuaian bacaan dengan konteks domain, meliputi rentang
berat yang sah untuk suatu produk, toleransi penimbangan, dan selisih terhadap
riwayat transaksi sebelumnya. Konfigurasi produk disimpan pada basis data dan
menjadi rujukan keputusan.

\subsection{Lapis 3: Verifikasi Temporal}

Lapisan ini membandingkan hasil pembacaan pada beberapa bingkai dari pembacaan
yang sama. Bacaan yang benar cenderung stabil, sedangkan kesalahan cenderung
berubah-ubah. Gagasan penyaringan antarbingkai telah dipakai untuk stabilisasi
keluaran \parencite{yenneti2023lightweight}, tetapi belum dimanfaatkan sebagai
dasar keputusan menerima atau menolak.

\subsection{Lapis 4: Verifikasi Statistik}

Lapisan ini memakai nilai keyakinan yang telah dikalibrasi. Kalibrasi dilakukan
pada bagian data yang terpisah dari data penilaian, sehingga ambang keputusan
tidak ditentukan dari data yang dipakai menilai kinerja.

\section{Rancangan Sistem}

Keputusan setiap lapisan digabungkan menjadi tiga kemungkinan keluaran, yaitu
terima, tolak, atau minta konfirmasi. Keluaran terima dan konfirmasi disimpan ke
basis data transaksi beserta alasannya, sedangkan keluaran tolak dicatat pada log
penolakan. Semua keputusan meninggalkan rekaman sehingga dapat diaudit dan
menjadi bahan perhitungan laju kesalahan senyap serta tingkat koreksi operator.
Sistem dirancang berjalan tanpa koneksi jaringan dan melakukan sinkronisasi
ketika jaringan tersedia.

\section{Rancangan Eksperimen}

Eksperimen disusun berjenjang agar kontribusi setiap tahap dapat dipisahkan.
Rincian disajikan pada Tabel~\ref{tab:eksperimen}.

\begin{table}[H]
\caption{Rancangan eksperimen}
\label{tab:eksperimen}
\begin{tabularx}{\textwidth}{@{}p{0.8cm} p{5.2cm} X@{}}
\toprule
\textbf{Kode} & \textbf{Konfigurasi} & \textbf{Pertanyaan yang dijawab} \\
\midrule
E1 & PP-OCRv6 \textit{tiny}, \textit{small}, dan \textit{medium} tanpa adaptasi & Sejauh mana kemampuan tanpa adaptasi? \\
E2 & Penggantian kepala keluaran menjadi 12 kelas & Tingkat adaptasi minimum yang memadai? \\
E3 & Ditambah verifikasi struktural & Berapa kesalahan senyap yang tertangkap? \\
E4 & Ditambah verifikasi fisik & Tambahan manfaat dari konteks produk? \\
E5 & Ditambah verifikasi temporal dan kalibrasi & Di mana titik seimbang antara \textit{coverage} dan akurasi? \\
E6 & Sistem penuh dibandingkan kondisi awal & Dampak terhadap tingkat koreksi dan waktu per transaksi \\
\bottomrule
\end{tabularx}
\end{table}

\section{Metrik Evaluasi}

Metrik yang dipakai diuraikan sebagai berikut.
\begin{poin}
  \item \textit{Full-string recognition} (FSR), yaitu persentase bacaan yang
        benar secara utuh.
  \item \textit{Character error rate} (CER), yaitu laju kesalahan tingkat
        karakter.
  \item Akurasi posisi desimal, yaitu persentase titik desimal pada posisi yang
        benar.
  \item Laju kesalahan senyap, yaitu persentase bacaan salah yang diterima
        sistem tanpa peringatan. Metrik ini merupakan metrik utama.
  \item \textit{Coverage}, yaitu persentase bacaan yang diterima otomatis.
  \item \textit{Selective accuracy}, yaitu akurasi pada bacaan yang diterima.
  \item AURC, yaitu luas daerah di bawah kurva \textit{risk--coverage}.
  \item Latensi p50 dan p95, ukuran model, serta tingkat koreksi operator.
\end{poin}

Setiap persentase yang dilaporkan disertai jumlah sampel dan interval
kepercayaan 95 persen. Perbandingan latensi disertai keterangan ukuran model dan
perangkat yang dipakai.

\section{Jadwal Penelitian}

Pelaksanaan penelitian direncanakan dalam sepuluh minggu kerja sebagai berikut.
\begin{poin}
  \item Minggu 1: penetapan spesifikasi data, metrik, dan ambang keberhasilan.
  \item Minggu 2: pembangunan pembangkit data sintetis dan penyusunan data uji
        terkontrol.
  \item Minggu 3: pengukuran kemampuan awal PP-OCRv6 untuk eksperimen E1.
  \item Minggu 4: pengambilan data uji nyata, penguncian data uji, dan
        pelaksanaan E2.
  \item Minggu 5: implementasi dan pengukuran verifikasi struktural (E3).
  \item Minggu 6: implementasi dan pengukuran verifikasi fisik (E4).
  \item Minggu 7: verifikasi temporal, kalibrasi, dan penyusunan kurva
        \textit{risk--coverage} (E5).
  \item Minggu 8: perangkaian sistem penuh dan pengujian fungsional (E6).
  \item Minggu 9: analisis kesalahan dan penyusunan Bab IV.
  \item Minggu 10: penyusunan Bab V dan konsolidasi naskah.
\end{poin}


% ---------------- Daftar pustaka -------------------------------------
\clearpage
\renewcommand{\thepage}{\arabic{page}}
\setcounter{page}{1}
\bibliography{references}
\addcontentsline{toc}{chapter}{DAFTAR PUSTAKA}

\end{document}
